\documentclass[a4paper,10pt]{article}
\newcommand\subdir{/home/koen/LaTeX-setup/}
\input{\subdir/LaTeX-files/imports.tex}
\input{\subdir/LaTeX-files/master-internship.tex}
\usepackage{\subdir/LaTeX-files/aastex}
\usepackage{natbib}
\bibliographystyle{\subdir/LaTeX-files/astroadstitle.bst}

%Set the title, Author and date
\title{Notes Week 7}
\author{Koen Essers}
\tutorial{Master Internship}
\date{\today}

%Set the header style
\header{\tutorialstring}

%Begin the document
\begin{document}
\setuplayout
\section{Results TPAGB Grid Week 6}

Presumably due to the very low \lstinline[style=output, breaklines=true]|envelope_mass_limit| of \lstinline[style=output, breaklines=true]|0.0001|, many of the TPAGB mass transfer models took very long to run.

Since the evolution from the end of the TPAGB phase to a white dwarf is not important for mass transfer, this extra computation is unnecessary.
I ended the pending jobs from last week and let the remaining models finish\sidenote{
Either until the star is a white dwarf, or the 2 days wall time had exceeded.}. This resulted in a total of 30 models.

Although 30 is not an unmanageable amount to manually load, I think, also with the future in mind, I could really use a nice \lstinline[style=output, breaklines=true]|MesaGrid| class to handle the loading and retrieving of mesa models in grids.

\subsection{MesaGrid Class}

The \lstinline[style=output, breaklines=true]|MesaGrid| class has a few core features.
\begin{citemize}
    \item Loading all, or specific \lstinline[style=output, breaklines=true]|binary| and \lstinline[style=output, breaklines=true]|star| \lstinline[style=output, breaklines=true]|history.data| files.
    \item Sorting \lstinline[style=output, breaklines=true]|history| data.
    \item Iteratively retrieving all models.
    \item Iteratively retrieving models with equal \(q\) or \(R_\text{RL}\).
    \item Retrieving a specific model.
    \item Accessing reference single star \lstinline[style=output, breaklines=true]|history| data
    \item Aligning binary \lstinline[style=output, breaklines=true]|history| data with single star \lstinline[style=output, breaklines=true]|history| data.
\end{citemize}

To achieve this, I use two classes. A \lstinline[style=output, breaklines=true]|MesaModel| class, used to store the important data and properties for the mesa models and the \lstinline[style=output, breaklines=true]|MesaGrid| class.

The \lstinline[style=output, breaklines=true]|MesaModel| class is very simple and is only used to conveniently access model data.
\begin{minted}[fontsize=\small]{py}
class MesaModel:
    def __init__(self, binary_path=None, tpagb_path=None, initial_age=None):

        self.binary = mr.MesaData(str(binary_path)) if binary_path else None
        self.star = mr.MesaData(str(tpagb_path)) if tpagb_path else None
        self.initial_age = initial_age
        self.age = self.star.star_age if tpagb_path else None
        self.env_mass = (
            self.star.star_mass - self.star.he_core_mass if tpagb_path else None
        )

    def __str__(self):
        if self.star:
            return f"MESA model with R = {self.star.R[0]:.2f},\
            R_RL = {self.star.rl_1[0]:.2f} and \
            q = {(self.star.star_2_mass[0] / self.star.star_1_mass[1]):.3f}"

        return f"Empty MESA model"

    def set_initial_age(self, age):
        self.initial_age = age
        self.age = self.star.star_age + self.initial_age

\end{minted}

The \lstinline[style=output, breaklines=true]|MesaGrid| class is much longer, so I will only show the most important parts:

\marginfignote{4}{Set the grid directory and filters. Call internal functions to load the models and references.}
\marginfignote{20}{Load the single star reference histories and save them.}
\marginfignote{25}{Find the reference star age at the start of the TPAGB.}
\marginfignote{29}{Normalises input of \lstinline[style=output, breaklines=true]|R1| and \lstinline[style=output, breaklines=true]|q| into a set.}
\marginfignote{35}{Returns whether the requested model passes the user given filter.}
\marginfignote{41}{Reads \lstinline[style=output, breaklines=true]|R_1| and \lstinline[style=output, breaklines=true]|q| from the model run directory and returns these values.}
\marginfignote{47}{For all models, parses the directory name, checks if the model passes the user defined filter and reads the binary and single star history. It then computes the initial age of the model and saves it to a last of all models. It also adds the parsed info into the sets of \lstinline[style=output, breaklines=true]|q| and \lstinline[style=output, breaklines=true]|R_1|.}
\marginfignote{80}{Finds the initial age of the model relative to the TPAGB by comparing the first entry of the history data with the reference single star history.}
\marginfignote{86}{Sorts the keys so they are ordered first by ascending Roche lobe radius, then by ascending \(q\).}
\marginfignote{90}{Allows easy access of specific model.}
\marginfignote{94}{Also access specific model, no error if the model does not exist.}
\marginfignote{98}{Iterate in sorted order: R1 ascending, then q ascending}
\marginfignote{104}{Return all models with given R1 sorted by q}
\marginfignote{114}{Return all models with given mass ratio
        sorted by R1.}
\begin{minted}[fontsize=\small,breaklines]{py}
class MesaGrid:
    def __init__(self, grid_dir, R1=None, q=None):
        ...
        self.R1_filter = self._normalize_filter(R1)
        self.q_filter = self._normalize_filter(q)

        self._get_tpagb_age()
        self._load_references()

        self._load_models()
        self._finalize_grid()

        
    def _load_references(self):
        ...
        
    def _get_tpagb_age(self):
        ...

    def _normalize_filter(self, val):
        ...
        return {filter}

    def _passes_filter(self, R1, q):
        ...
        return bool

    def _parse_name(self, name):
        ...
        return R1, q

    def _load_models(self):
        ...
        for run_dir in run_dirs:
            ...
            parsed = self._parse_name(run_dir.name)
            ...
            R1, q = parsed

            if not self._passes_filter(R1, q):
                continue

            ...
            model = MesaModel(
                binary_path,
                tpagb_path,
            )

            age = self._find_initial_age(model)
            model.set_initial_age(age)

            self.models[(R1, q)] = model
            self.R1_vals.add(R1)
            self.q_vals.add(q)

    def _find_initial_age(self, model):
        ...
        return model_age - self.tpagb_age

    def _finalize_grid(self):
        ...

    def __getitem__(self, key):
        return self.models[key]

    def get(self, R1, q):
        return self.models.get((R1, q))

    def iter_models(self):
        for R1, q in self.sorted_keys:
            yield R1, q, self.models[(R1, q)]

    def get_R1_slice(self, R1):
        return [
            (q, self.models[(R1, q)]) 
            for q in self.q_vals
            if (R1, q) in self.models
        ]

    def get_q_slice(self, q):
        return [
            (R1, self.models[(R1, q)])
            for R1 in self.R1_vals 
            if (R1, q) in self.models
        ]
\end{minted}

\newsubsection{Demo}
Below, I show some use cases:
\marginfignote{3}{Loads all models with \(q=0.5\)}
\begin{minted}[fontsize=\small,breaklines]{py}
grid = MesaGrid(".../grid-dir/", q=0.500)
print(grid.models)
\end{minted}
\begin{minted}[fontsize=\small]{text}
{(225.0, 0.5): <scripts.general_utils.mesa_grid.MesaModel object at 0x7ff1c3aa27b0>,
 (150.0, 0.5): <scripts.general_utils.mesa_grid.MesaModel object at 0x7ff1d6e10410>,
 (375.0, 0.5): <scripts.general_utils.mesa_grid.MesaModel object at 0x7ff1d6e11090>,
 (300.0, 0.5): <scripts.general_utils.mesa_grid.MesaModel object at 0x7ff1d7f31350>}
\end{minted}
We can iterate over these models using the \lstinline[style=output, breaklines=true]|grid.iter_models()| function:
\begin{minted}[fontsize=\small,breaklines]{py}
for R1, q, model in grid.iter_models():
    print(f"{R1 = }, {q = }, {model}")
\end{minted}
\begin{minted}[fontsize=\small]{text}
R1 = 150.0, q = 0.5, MESA model with R = 116.62, R_RL = 150.00 and q = 0.500 
R1 = 225.0, q = 0.5, MESA model with R = 116.62, R_RL = 225.00 and q = 0.500
R1 = 300.0, q = 0.5, MESA model with R = 213.79, R_RL = 300.00 and q = 0.500
R1 = 375.0, q = 0.5, MESA model with R = 272.64, R_RL = 375.00 and q = 0.500
\end{minted}

Loading a 2d grid is also easy:
\begin{minted}[fontsize=\small,breaklines]{py}
grid = MesaGrid(
    ".../grid-dir/", 
    q=[0.125,0.250,0.375], 
    R1=[150,225,300])
for R1, q, model in grid.iter_models():
    print(f"{R1 = }, {q = :.3f}, {model}")
\end{minted}
\begin{minted}[fontsize=\small]{text}
R1 = 150.0, q = 0.125, MESA model with R = 116.62, R_RL = 150.00 and q = 0.125
R1 = 150.0, q = 0.250, MESA model with R = 116.62, R_RL = 150.00 and q = 0.250
R1 = 150.0, q = 0.375, MESA model with R = 116.62, R_RL = 150.00 and q = 0.375
R1 = 225.0, q = 0.125, MESA model with R = 116.62, R_RL = 225.00 and q = 0.125
R1 = 225.0, q = 0.250, MESA model with R = 116.62, R_RL = 225.00 and q = 0.250
R1 = 225.0, q = 0.375, MESA model with R = 116.62, R_RL = 225.00 and q = 0.375
R1 = 300.0, q = 0.125, MESA model with R = 213.79, R_RL = 300.00 and q = 0.125
R1 = 300.0, q = 0.250, MESA model with R = 213.79, R_RL = 300.00 and q = 0.250
R1 = 300.0, q = 0.375, MESA model with R = 213.79, R_RL = 300.00 and q = 0.375
\end{minted}

The grid can be sliced using the \lstinline[style=output, breaklines=true]|get_q_slice| and \lstinline[style=output, breaklines=true]|get_R1_slice| functions:
\begin{minted}[fontsize=\small,breaklines]{py}
print("slicing q")
for R1, model in grid.get_q_slice(q=0.125):
    print(model)

print("\nslicing R1")
for R1, model in grid.get_R1_slice(R1=225):
    print(model)
\end{minted}
\begin{minted}[fontsize=\small]{text}
slicing q
MESA model with R = 116.62, R_RL = 150.00 and q = 0.125
MESA model with R = 116.62, R_RL = 225.00 and q = 0.125
MESA model with R = 213.79, R_RL = 300.00 and q = 0.125

slicing R1
MESA model with R = 116.62, R_RL = 225.00 and q = 0.125
MESA model with R = 116.62, R_RL = 225.00 and q = 0.250
MESA model with R = 116.62, R_RL = 225.00 and q = 0.375
\end{minted}

Using the \lstinline[style=output, breaklines=true]|MesaGrid| class, we can select the reference single star history. It is also possible to select a single model from the grid:
\begin{minted}[fontsize=\small,breaklines]{py}
model = grid.get(R1=300, q=0.875)
plt.plot(grid.ref_tpagb.star_age, grid.ref_tpagb.R)
plt.plot(model.age, model.star.R)
plt.show()
\end{minted}
\begin{figure}[ht]
    \centering
    \incplot{w7-compare-with-ref}
\end{figure}
\begin{minted}[fontsize=\small,breaklines]{py}
model = grid.get(R1=300, q=0.875)
plt.plot(model.env_mass, model.star.R)
plt.plot(model.env_mass, model.star.rl_1)
plt.gca().invert_xaxis()
plt.show()
\end{minted}
\begin{figure}[ht]
    \centering
    \incplot{w7-show-mass}
\end{figure}

\subsection{Results}

Since I ended last week's run early, I only had a complete set of models for \(R_{RL} = 300R_\odot \). Therefore, I will plot some (hopefully) interesting things using these models.

Below, I plot the relative Roche lobe sizes as a function of envelope mass for different \(q\) and a starting Roche lobe radius of \(R_\text{RL} = 300 R_\odot \).

\begin{figure}[ht]
    \centering
    \incplot{relativerocheloberadius}
\end{figure}
I think this nicely shows how import the mass ratio is for the resulting maximum \(R_\text{star} / R_\text{RL}\).

\nextpage

In a similar way, using the analytic approximation for the outer Roche lobe radius from \citet{temmink2022}, the effects of \(q\) on the ratio \(R_\text{star} / R_\text{ROL}\) can be plotted. Here, I shade the region above this value, the MESA simulations are not reliable.
\begin{figure}[ht]
    \centering
    \incplot{relativeouterrocheloberadius}
\end{figure}

Similarly, the radius of the star and the radius of the orbit can be compared. When the ratio exceeds 2, the system must be a contact binary, even for the smallest accretor. 

\begin{figure}[ht]
    \centering
    \incplot{R300period}
\end{figure}

Then, the ratio of actual mass loss over the quasi adiabatic limit. Here a value of 1 should mean mass loss cannot be stable.

\begin{figure}[ht]
    \centering
    \incplot{w7-mdot}
\end{figure}

Clearly, all models stay well below 1.


\newsection{Overshooting}

I compare the single star evolution of stars with different masses for 3 different overshooting prescriptions:
\begin{enumd}
    \item  Overshooting from \citet{rees2024b}
    \item  Overshooting from \citet{temmink2022}
    \item  Overshooting from \citet{temmink2022} with different \(f_\text{ov,0}\).
\end{enumd}

Below, I note the exact inlist settings for all three cases.
\subsection{\texorpdfstring{\citet{rees2024b}}{Rees et al. (2024)}}
\begin{minted}[fontsize=\small,breaklines]{fortran}
! inlist_MS
&controls
    ...
	overshoot_scheme(1) = 'step' !jermyn et al 2022
	overshoot_zone_type(1) = 'any'
	overshoot_zone_loc(1) = 'core'
	overshoot_bdy_loc(1) = 'any'
	overshoot_f(1) = 0.129
	overshoot_f0(1) = 0.0129
	overshoot_D0(1) = 0d0
	overshoot_Delta0(1) = 1d0
	overshoot_scheme(2) = 'exponential'
	overshoot_zone_type(2) = 'nonburn'
	overshoot_zone_loc(2) = 'shell'
	overshoot_bdy_loc(2) = 'any'
	overshoot_f(2) = 0.0174
	overshoot_f0(2) = 0.00174
	overshoot_D0(2) = 0d0
	overshoot_Delta0(2) = 1d0
    ...
\end{minted}

\begin{minted}[fontsize=\small,breaklines]{fortran}
! inlist_GB
&controls
    ...
	overshoot_scheme(1) = 'exponential'
	overshoot_zone_type(1) = 'nonburn'
	overshoot_zone_loc(1) = 'shell'
	overshoot_bdy_loc(1) = 'any'
	overshoot_f(1) = 0.0174
	overshoot_f0(1) = 0.00174
	overshoot_D0(1) = 0d0
	overshoot_Delta0(1) = 1d0
    ...
\end{minted}

\begin{minted}[fontsize=\small,breaklines]{fortran}
! inlist_CHeB
&controls
    ...
	predictive_mix(1) = .true.
	predictive_zone_type(1) = 'any'
	predictive_zone_loc(1) = 'core'
	predictive_bdy_loc(1) = 'any'
	predictive_superad_thresh(1) = 0.005
	predictive_avoid_reversal = 'he4'

	overshoot_scheme(2) = 'exponential'
	overshoot_zone_type(2) = 'nonburn'
	overshoot_zone_loc(2) = 'shell'
	overshoot_bdy_loc(2) = 'any'
	overshoot_f(2) = 0.0174
	overshoot_f0(2) = 0.00174
	overshoot_D0(2) = 0d0
	overshoot_Delta0(2) = 1d0
   ...
\end{minted}

\begin{minted}[fontsize=\small,breaklines]{fortran}
! inlist_EAGB
&controls
    ...
	overshoot_scheme(1) = 'exponential' ! IPCZ overshooting controls
	overshoot_zone_type(1) = 'burn_He'
	overshoot_zone_loc(1) = 'shell'
	overshoot_bdy_loc(1) = 'bottom'
	overshoot_f(1) = 0.008 ! Herwig 2008
	overshoot_f0(1) = 0.0008
	overshoot_scheme(2) = 'exponential' ! envelope overshooting
	overshoot_zone_type(2) = 'nonburn'
	overshoot_zone_loc(2) = 'shell'
	overshoot_bdy_loc(2) = 'any'
	overshoot_f(2) = 0.0174
	overshoot_f0(2) = 0.00174
	overshoot_scheme(3) = 'exponential' ! carbon burning
	overshoot_zone_type(3) = 'any'
	overshoot_zone_loc(3) = 'any'
	overshoot_bdy_loc(3) = 'bottom'
	overshoot_f(3) = 0.005
	overshoot_f0(3) = 0.0005
  ...
\end{minted}

\subsection{\texorpdfstring{\citet{temmink2022}}{Temmink et al. (2023)}}
\begin{minted}[fontsize=\small,breaklines]{fortran}
! inlist_common
&controls
    ...
	overshoot_scheme(1) = 'exponential'
	overshoot_zone_type(1) = 'any'
	overshoot_zone_loc(1) = 'core'
	overshoot_bdy_loc(1) = 'top'
	overshoot_f(1) = 0.016d0
	overshoot_f0(1) = 0.08d0

	overshoot_scheme(2) = 'exponential'
	overshoot_zone_type(2) = 'any'
	overshoot_zone_loc(2) = 'shell'
	overshoot_bdy_loc(2) = 'any'
	overshoot_f(2) = 0.0174d0
	overshoot_f0(2) = 0.0087d0
    ...
\end{minted}

\subsection{\texorpdfstring{\citet{temmink2022}}{Temmink et al. (2023)} other \texorpdfstring{\(f_\text{ov,0}\)}{f\_ov,0}}
\begin{minted}[fontsize=\small,breaklines]{fortran}
! inlist_common
&controls
    ...
	overshoot_scheme(1) = 'exponential'
	overshoot_zone_type(1) = 'any'
	overshoot_zone_loc(1) = 'core'
	overshoot_bdy_loc(1) = 'top'
	overshoot_f(1) = 0.016d0
	overshoot_f0(1) = 0.008d0

	overshoot_scheme(2) = 'exponential'
	overshoot_zone_type(2) = 'any'
	overshoot_zone_loc(2) = 'shell'
	overshoot_bdy_loc(2) = 'any'
	overshoot_f(2) = 0.0174d0
	overshoot_f0(2) = 0.0087d0
   ...
\end{minted}

\newsubsection{Results}


\subsubsection{HR-diagrams}
Below, I show Hertzsprung Russel diagrams for the pre-TPAGB phases. For 7 different masses, the overshooting options listed above are shown.

I also compare these tracks with MIST \citep{choi2016} EEP tracks.

\begin{figure}[ht]
    \centering
    \incplot{hr-overshooting}
\end{figure}

Now, it is quite difficult to see real differences in these small plots. Therefore, I plot the phases separately as well.

\begin{figure}[ht]
    \centering
    \incplot{hr-overshooting-MS}
\end{figure}

\begin{figure}[ht]
    \centering
    \incplot{hr-overshooting-GB}
\end{figure}

\nextpage

\begin{figure}[ht]
    \centering
    \incplot{hr-overshooting-CHeB}
\end{figure}

The tracks in the HR-diagrams for stars using the overshooting options from \citet{rees2024b} and \citet{temmink2022} are very similar. When using the other \(f _\text{ov,0}\), the results are however significantly different and closer to MIST EEPs. In general, these stars have higher luminosities.

For the MIST \citep{choi2016} EEP tracks, I used a set of non-rotating [Fe/H] = 0 models. This corresponds with an initial \(Z=0.01429\), slightly higher than the \(Z=0.014\) used for my MESA models. The MIST models also use Ferguson low-T opacities as well as a slightly different mixing length parameter \(\alpha _\text{MLT} = 1.82\) instead of \(\alpha _\text{MLT} = 1.931\) from \citet{cinquegrana2022}. Maybe these differences could account for the differences visible between the MIST and \citet{temmink2022} models with alternative \(f _\text{ov, 0}\).


\subsubsection{Kippenhahn diagrams}

Below, I show Kippenhahn diagrams of the evolutionary stages for all stars.
For completeness, I plot all of them, but really, the different masses are not very different\sidenote{It took me quite some time to figure out that the \lstinline[style=output, breaklines=true]|mkipp| code that I was using, used MESAs old definitions of mixing regions, while my models used the new definitions. I therefore had to change the \lstinline[style=output, breaklines=true]|mkipp| code to use the newer definitions.}.

The gray regions are convective regions in the star. The red regions are overshooting regions. The blue "contours" show where nuclear burning is happening. The black line is the He-core mass and the blue line is the CO-core mass.

The last row compares the overshooting regions. The blue regions are the overshooting regions from \citet{rees2024b}. The orange regions are the overshooting regions from \citet{temmink2022} and the green regions are the overshooting regions from \citet{temmink2022} with a different value of \(f _\text{ov,0}\).

\begin{figure}[ht]
    \centering
    \incplot{kippenhahn2}
\end{figure}



\begin{figure}[ht]
    \centering
    \incplot{kippenhahn3}
\end{figure}

\nextpage

\begin{figure}[ht]
    \centering
    \incplot{kippenhahn4}
\end{figure}

\begin{figure}[ht]
    \centering
    \incplot{kippenhahn5}
\end{figure}

\begin{figure}[ht]
    \centering
    \incplot{kippenhahn6}
\end{figure}

\nextpage

\begin{figure}[ht]
    \centering
    \incplot{kippenhahn7}
\end{figure}

\begin{figure}[ht]
    \centering
    \incplot{kippenhahn8}
\end{figure}

The core mass at the start of the TPAGB is indistiguishable for the different models. Does this mean it is not very important which prescription is used?

\newsection{Opacity tables}

To see if using the \citet{ferguson1905} opacity table might improve stability during episodes of fast mass transfer during a thermal pulse, I compare binary models with donor TPAGB stars using the \citet{ferguson1905} opacity tables, and the \citet{marigo2009} opacity tables.

I set up a grid for 24 models for both stars where for every initial Roche lobe radius, I simulate the binary for \(q=0.25, 0.5, 0.75\) and \(1\).
For the radii, I chose \(R_\text{RL} (R_\odot )= 150, 240, 330, 420, 510\) and \(600\).

Unfortunately, it quite quickly became clear that both models have trouble with stability at low \(q\) and low \(R_\text{RL}\).

Below, I plot which models failed. The thick lines correspond to \AE SOPUS models, while the thin lines correspond to Ferguson models.

\begin{figure}[ht]
    \marginfignote{4}{I think it is interesting that the problem only occurs for small radii. Maybe this has to do with the ratio of star radius versus Roche lobe radius.}
    \centering
    \incplot{w7-compare-opacity}
\end{figure}

To test my hypothesis, I plot the ratio of star radius over Roche lobe radius for \(q=0.25\) and different starting radii for the Ferguson models.

\begin{figure}[ht]
    \marginfignote{4}{There does seem to be a correlation between the initial Roche lobe radius and the ratio \(R_\text{star} / R_\text{RL}\).}
    \centering
    \incplot{w7-compare-q}
\end{figure}

I also plot elapsed time for all Ferguson models.

\begin{figure}[ht]
    \marginfignote{4}{Two things here are interesting: 
    \begin{citemize}
    \item Generally, models with a larger starting Roche lobe radius take longer to start losing mass. 
    \item  The last part of the evolution can take a significant portion of the computing time, even when ending at \(0.01M_\odot \).
\end{citemize}
    }
    \centering
    \incplot{w7-all-compare-opacity}
\end{figure}

\nextpage

I also check if the temperature gradient problems are present for both failing \citet{ferguson1905} models, as well as failing \AE SOPUS \citep{marigo2009} models. Unfortunately, this is the case, as can be seen below:

Ferguson \(R_\text{RL} = 150\), \(q = 0.25\):
\begin{figure}[ht]
    \centering
    \incplot{w7-problem-gradTgradR}
\end{figure}

\AE SOPUS \(R_\text{RL} = 150\), \(q = 0.25\)
\begin{figure}[ht]
    \centering
    \incplot{w7-problem-gradTgradR-aeso}
\end{figure}

I also thought the \(R_\text{RL} = 600R_\odot \) model looked interesting in the plots above. Below I plot the evolution of the Roche lobe:

\begin{figure}[ht]
    \marginfignote{4}{I think this figure quite clearly shows that in parts where wind mass loss is important, the Roche lobe radius will increase significantly as a result. In extreme cases, this can mean that the star just barely avoids filling its Roche lobe.}
    \centering
    \incplot{w7-large-R}
\end{figure}

\newsection{Close \texorpdfstring{\(R_\text{RL}\)}{R\_RL} Spacing}

I also simulated a grid of models with \(q=0.8\) and \(R_\text{RL}\) ranging from \(506R_\odot \) to \(550R_\odot \).

These models just finished this morning and I only had time to make a basic plot without looking too much into the models too much.
\begin{figure}[ht]
    \centering
    \incplot{w7-end}
\end{figure}


\bibliography{master.bib}
\end{document}
